
\SGARefTarget{sga3:xi:section:4}
\section*{4. The Scheme of Multiplicative-Type Subgroups\\
of a Smooth Affine Group}
\addcontentsline{toc}{section}{4. The Scheme of Multiplicative-Type
Subgroups of a Smooth Affine Group}

The principal result of the present Exposé is \SGARefLink{sga3:xi:theorem:4:1}{the following}.

\medskip
\noindent\SGARefTarget{sga3:xi:theorem:4:1}\textbf{Theorem 4.1.} ---
\textit{Let \(S\) be a prescheme, let \(G\) be an \(S\)-group
prescheme smooth and affine over \(S\), and let}
\[
  F\colon(\mathbf{Sch})^\circ/S\longrightarrow(\mathbf{Ens})
\]
\textit{be the functor defined by}
\SGARefTarget{sga3:xi:display:4:1:functor}
\[
  F(T)=\{\text{multiplicative-type subgroups of \(G_T\)}\}.
\]
\textit{Then \(F\) is representable, smooth, and separated over
\(S\).}

Let us record at once \SGARefLink{sga3:xi:corollary:4:2}{the following variant}.

\medskip
\noindent\SGARefTarget{sga3:xi:corollary:4:2}\textbf{Corollary 4.2.} ---
\textit{Let \(G\) and \(H\) be two \(S\)-group preschemes, with
\(G\) smooth and affine over \(S\), and \(H\) of multiplicative type
and of finite type.  Then
\(\operatorname{Hom}_{S\text{-}\mathrm{gr}}(H,G)\) is representable,
smooth, and separated over \(S\).}

Indeed, when \(H\) is smooth over \(S\), so is \(H\times_S G\), and
one may apply \SGARefLink{sga3:xi:theorem:4:1}{4.1} to the latter.
By considering graph subgroups,
\(\operatorname{Hom}_{S\text{-}\mathrm{gr}}(H,G)\) becomes a
subfunctor \(F'\) of the functor \(F\) of multiplicative-type
subgroups of \(H\times_S G\): namely, \(F'(T)\) is the set of
multiplicative-type subgroups \(K\) of
\[
  (H\times_S G)_T=H_T\times_TG_T
\]
such that the homomorphism induced by the first projection,
\SGARefTarget{sga3:xi:display:4:2:first-projection}
\[
  K\longrightarrow H_T,
\]
is an isomorphism.  By \SGARefLink{sga3:ix:corollary:2:9}{Exposé IX, 2.9}, the canonical morphism
\(F'\to F\) is an open immersion.  Since \(F\) is representable,
the same is true of \(F'\), which is represented by an open
subprescheme of \(F\); and since \(F\) is separated and smooth over
\(S\), so is \(F'\).

When \(H\) is finite over \(S\), it is enough to apply
\SGARefLink{sga3:xi:proposition:3:12:item:b}{3.12(b)} to obtain the
representability of
\(\operatorname{Hom}_{S\text{-}\mathrm{gr}}(H,G)\).  When
\(H=H_1\times_SH_2\), with \(H_1\) smooth over \(S\) and \(H_2\)
finite over \(S\), then
\(\operatorname{Hom}_{T\text{-}\mathrm{gr}}(H_T,G_T)\) is identified
with the subset of
\SGARefTarget{sga3:xi:display:4:2:product-hom}
\[
 \operatorname{Hom}_{T\text{-}\mathrm{gr}}\bigl((H_1)_T,G_T\bigr)
 \times
 \operatorname{Hom}_{T\text{-}\mathrm{gr}}\bigl((H_2)_T,G_T\bigr)
\]
formed by the pairs \((u_1,u_2)\) such that
\(u_1\mathbin{\cdot}u_2=u_2\mathbin{\cdot}u_1\).  It follows that
\(\operatorname{Hom}_{S\text{-}\mathrm{gr}}(H,G)\) is represented
by a closed subprescheme of
\SGARefTarget{sga3:xi:display:4:2:representing-product}
\[
 \operatorname{Hom}_{S\text{-}\mathrm{gr}}(H_1,G)
 \times_S
 \operatorname{Hom}_{S\text{-}\mathrm{gr}}(H_2,G)=X.
\]
Indeed, apply \SGARefLink{sga3:viii:examples:6:5:item:b:01}{Exposé VIII, 6.5(b)}, with \(Y=H\), \(Z=G\), and with
\(q_1\) and \(q_2\) defined respectively by
\[
  (u_1,u_2)\longmapsto u_1\mathbin{\cdot}u_2
  \qquad\text{and}\qquad
  (u_1,u_2)\longmapsto u_2\mathbin{\cdot}u_1.
\]

In the general case, the question is local on \(S\) for the Zariski
topology, so we may suppose that \(S\) is affine and that \(H\) has
constant type over \(S\).  Since \(H\) is quasi-isotrivial
(\SGARefLink{sga3:x:corollary:4-5}{Exposé X, 4.5}), there is a surjective étale morphism \(S'\to S\),
with \(S'\) affine, that splits \(H\), that is, such that
\(H'=H_{S'}\) is diagonalizable.  \SGARefLink{sga3:xi:corollary:4:2}{The preceding result} then applies,
because a diagonalizable group is the product of a diagonalizable
torus and a finite diagonalizable group.  It remains to see that the
descent datum obtained on the \(S'\)-prescheme
\[
  \operatorname{Hom}_{S'\text{-}\mathrm{gr}}(H',G')=X'
\]
is effective.  This follows from the argument of \SGARefLink{sga3:x:lemma:5-4}{Exposé X, 5.4},
noting that, there, the hypothesis that \(X'\to S'\) be separated
and locally quasi-finite was used only to ensure that every open
part of \(X'\) quasi-compact over \(S'\) was quasi-affine over
\(S'\).  That property still holds here, as follows readily from
the fact that it holds for the functor \(F\) of
\SGARefLink{sga3:xi:theorem:4:1}{4.1} (as will be seen during the
proof of \SGARefLink{sga3:xi:theorem:4:1}{4.1}).  This argument---or
any similar variant of this small dévissage---establishes the
representability of
\(\operatorname{Hom}_{S\text{-}\mathrm{gr}}(H,G)\), and at the same
time shows that it is separated over \(S\).  It is locally of finite
presentation over \(S\), as follows, for example, from the fact
already noted in \SGARefLink{sga3:xi:remark:3:2}{3.2} that this
functor ``commutes with filtered direct limits of rings.''  Finally,
the functor is formally smooth over \(S\) by
\SGARefLink{sga3:xi:proposition:2:1}{2.1}, and is therefore smooth
over \(S\).

\medskip
\noindent\SGARefTarget{sga3:xi:remark:4:3}\textbf{Remark 4.3.} ---
We have deduced \SGARefLink{sga3:xi:corollary:4:2}{4.2} from
\SGARefLink{sga3:xi:theorem:4:1}{4.1}; this is genuinely immediate
only when \(H\) is also smooth over \(S\).  To make the deduction
without contortions in the general case, one would need the
representability assertion in
\SGARefLink{sga3:xi:theorem:4:1}{4.1} without assuming that \(G\) is
smooth over \(S\), assuming only that it is affine and of finite
presentation over \(S\).  Of course, \(F\) would then no longer be
smooth over \(S\) in general.  There is little doubt that
\SGARefLink{sga3:xi:theorem:4:1}{4.1} remains true under these more
general hypotheses, but the proof appears likely to be more delicate,
because one can no longer invoke
\SGARefLink{sga3:xi:corollary:3:8}{3.8}.\footnote{This is indeed
proved for \(G\) flat and quasi-affine over \(S\), with connected
fibres, provided one restricts to the central subtori of \(G\)
(\SGARefLink{sga3:xv:theorem:8-8}{Exposé XV, 8.8}).}\footnote{Editor's note: when \(H\) is smooth (not
necessarily affine) over a normal locally noetherian base \(S\),
M.~Raynaud showed that the largest representable open subfunctor of
the functor of subtori of \(H\) is a disjoint sum of opens smooth and
affine over \(S\).  This is Theorem IX.9.26 in \emph{Faisceaux amples
sur les schémas en groupes et sur les espaces homogènes}, Lecture
Notes in Mathematics 119 (1970).}

Let us nevertheless point out that when \(G\) is a closed subgroup
of an affine smooth group \(G_0\) over \(S\), the functor \(F\) of
multiplicative-type subgroups of \(G\) is represented by a closed
subprescheme of the prescheme representing the analogous functor
\(F_0\) for \(G_0\), which is covered by
\SGARefLink{sga3:xi:theorem:4:1}{4.1}; this follows readily by
applying \SGARefLink{sga3:viii:theorem:6:4}{Exposé VIII, 6.4}.  This also raises the question whether an
affine \(S\)-group scheme \(G\), affine and of finite presentation
over \(S\), is isomorphic to a subgroup scheme of
\(\operatorname{GL}(n)_S\) for some suitable \(n\).  This is true
when \(S\) is the spectrum of a field (cf.\ Exposé
\(\mathrm{VI}_{\mathrm B}\), 11.11), but is unfortunately false in
general, even for tori; see
\SGARefLink{sga3:xi:remark:4:6}{4.6}.  Finally, one could also prove
\SGARefLink{sga3:xi:corollary:4:2}{4.2} directly by exactly the same
method as \SGARefLink{sga3:xi:theorem:4:1}{4.1}.

We now prove \SGARefLink{sga3:xi:theorem:4:1}{4.1}.  Since the
functor \(F\) is evidently local in nature, we may suppose that
\(S\) is affine, say \(S=\operatorname{Spec}(A)\), where \(A\) is a
ring.  Regarding \(A\) as the filtered direct limit of its subrings
of finite type over \(\mathbb Z\), and noting that \(G\) descends
from a smooth affine group over one such subring (EGA IV, 8), we are
reduced to the case in which \(S\) is noetherian.

For every integer \(n>0\), let \(T_n\) be the functor defined like
\(F\), but restricted to multiplicative-type subgroups \(H\) of
\(G_T\) such that \(n\cdot\operatorname{id}_H=0\), that is, such
that \({}_nH=H\).  Order the set \(I\) of positive integers by
divisibility.  When \(m\) is a multiple of \(n\), define
\SGARefTarget{sga3:xi:display:4:1:p-n-m}
\[
  p_{n,m}\colon T_m\longrightarrow T_n
\]
by
\[
  p_{n,m}(H)={}_nH=\operatorname{Ker}
  \bigl(n\cdot\operatorname{id}_H\bigr).
\]
In this way the \(T_n\) form an inverse system of functors over
\(S\).  By
\SGARefLink{sga3:xi:proposition:3:12:item:a}{3.12(a)}, the functors
\(T_n\) are representable, affine, and of finite type over \(S\).
Similarly define morphisms
\SGARefTarget{sga3:xi:display:4:1:u-n}
\[
  u_n\colon F\longrightarrow T_n
\]
by \(u_n(H)={}_nH\).  Thus one obtains a morphism
\SGARefTarget{sga3:xi:display:4:1:u-inverse-limit}
\[
  u\colon F\longrightarrow T=\varprojlim_n T_n,
\]
where the inverse limit is taken in the category of functors over
\(S\).  Since the \(T_n\) are representable and affine over \(S\),
so is \(T\): it is the spectrum of the filtered direct limit of the
quasi-coherent algebras on \(S\) defining the \(T_n\).  Of course,
\(T\) is not in general of finite type over \(S\).

We shall apply \SGARefLink{sga3:xi:proposition:3:7}{3.7}.  It is
enough to verify conditions \textit{\SGARefLink{sga3:xi:proposition:3:7:item:a}{(a)}}--\textit{\SGARefLink{sga3:xi:proposition:3:7:item:d}{(d)}} there; these
will imply that \(F\) is represented by a prescheme locally of finite
type over \(S\).  It then follows from
\SGARefLink{sga3:xi:proposition:2:1:bis}{2.1 bis} that \(F\) is even
smooth over \(S\).  Moreover, since \(F\) is a subfunctor of \(T\),
which is affine over \(S\), \(F\) is separated over \(S\): it is
separated over \(T\), and \(T\) is separated over \(S\).  We can
also prove at once the supplementary assertion used above: every
open \(U\subset F\) that is quasi-compact over \(S\) is quasi-affine
over \(S\).  This follows from
\SGARefLink{sga3:xi:corollary:3:11}{3.11} and the fact that the
\(T_n\) are affine over \(S\).

Let us now verify the conditions of
\SGARefLink{sga3:xi:proposition:3:7}{3.7}.

\smallskip
\noindent\SGARefTarget{sga3:xi:theorem:4:1:condition:a}\textit{(a)}
\(F\) is a sheaf for the faithfully flat quasi-compact topology, by
the descent theory of SGA~1, Exposé VIII, which applies here because
groups of multiplicative type over \(T\) are affine over
\(T\).\footnote{The printed source says ``affine over \(S\)'' here;
the descent assertion concerns groups over \(T\), so the required
base is \(T\).} It
commutes with filtered direct limits of rings by the standard
machinery of EGA IV, 8.  Let us show that it commutes with adic
inverse limits of local artinian rings.  For the analogous functor
considered in \SGARefLink{sga3:xi:corollary:4:2}{4.2}, this is
precisely \SGARefLink{sga3:ix:theorem:7:1}{Exposé IX, 7.1}, in the special case where \(A\) is a
complete noetherian local ring endowed with an ideal of definition
for its usual topology.  Notice that this is exactly where the
hypothesis that \(G\) be affine enters essentially.  In the present
case we are reduced to \SGARefLink{sga3:xi:lemma:4:4}{the following lemma}.

\medskip
\noindent\SGARefTarget{sga3:xi:lemma:4:4}\textbf{Lemma 4.4.} ---
\textit{Let \(A\) be a complete noetherian local ring with maximal
ideal \(\mathfrak m\), let \(G\) be a group scheme affine over
\(S=\operatorname{Spec}(A)\), and, for every integer \(n\geq0\), put}
\SGARefTarget{sga3:xi:display:4:4:adic-thickenings}
\[
  S_n=\operatorname{Spec}\bigl(A/\mathfrak m^{\,n+1}\bigr),
  \qquad
  G_n=G\times_SS_n.
\]
\textit{For every \(n\), let \(H_n\) be a multiplicative-type
subgroup of finite type of \(G_n\), and suppose that, whenever
\(m\geq n\), \(H_n\) is obtained from \(H_m\) by reduction.  Then
there exists a unique multiplicative-type subgroup \(H\) of \(G\)
whose reductions are the \(H_n\).}

By \SGARefLink{sga3:x:theorem:3:2}{Exposé X, 3.2}, there exists a multiplicative-type group \(H\) over
\(S\), necessarily of finite type and isotrivial and determined up
to unique isomorphism, together with an isomorphism
\[
  H\times_SS_0\simeq H_0.
\]
Here one uses the fact that \(H_0\) is isotrivial, since it is of
finite type over a field; see \SGARefLink{sga3:x:proposition:1:4}{Exposé X, 1.4}.  By \SGARefLink{sga3:x:proposition:2:1}{Exposé X, 2.1}, for
every \(n\) the isomorphism \(H\times_SS_0\simeq H_0\) lifts to a
unique isomorphism \(H\times_SS_n\simeq H_n\).  \SGARefLink{sga3:ix:theorem:7:1}{Exposé IX, 7.1},
already cited, now shows that the homomorphisms \(H_n\to G_n\) arise
from a unique homomorphism of \(S\)-groups
\SGARefTarget{sga3:xi:display:4:4:u}
\[
  u\colon H\longrightarrow G.
\]
By \SGARefLink{sga3:ix:corollary:6:6}{Exposé IX, 6.6}, \SGARefLink{sga3:xi:display:4:4:u}{this homomorphism} is a monomorphism because
\(u_0\colon H_0\to G_0\) is one.  This proves
\SGARefLink{sga3:xi:lemma:4:4}{4.4}.

\smallskip
\noindent\SGARefTarget{sga3:xi:theorem:4:1:condition:b}\textit{(b)}
The morphism \(u\colon F\to T\) is a monomorphism.  This follows from
the density theorem, \SGARefLink{sga3:ix:theorem:4:7}{Exposé IX, 4.7}, in the form of \SGARefLink{sga3:ix:corollary:4:8:item:b}{Corollary
4.8(b)}.  It is essential for the present application that this
result is available over a base that need not be noetherian.  Since
the functor \(T\) does not commute with filtered direct limits of
rings, one cannot reduce to the noetherian case a priori.

\smallskip
\noindent\SGARefTarget{sga3:xi:theorem:4:1:condition:b-prime}
\textit{(\(b'\))} The morphism \(u\colon F\to T\) is infinitesimally
étale; in other words, we have \SGARefLink{sga3:xi:lemma:4:5}{the following lemma}.

\medskip
\noindent\SGARefTarget{sga3:xi:lemma:4:5}\textbf{Lemma 4.5.} ---
\textit{Let \(A\) be an artinian local ring with residue field
\(k\), let \(S=\operatorname{Spec}(A)\), let
\(I\subset\operatorname{rad}(A)\) be an ideal, and put
\(S_0=\operatorname{Spec}(A/I)\).  Let \(G\) be an \(S\)-group
prescheme smooth over \(S\), and put \(G_0=G\times_SS_0\).  Suppose
given a multiplicative-type subgroup \(H_0\) of \(G_0\) and, for
every integer \(n>0\), a multiplicative-type subgroup \(H(n)\) of
\(G\), in such a way that}
\begin{enumerate}
\item[\SGARefTarget{sga3:xi:lemma:4:5:item:1}\(1^\circ\)]
\textit{for every multiple \(m\) of \(n\),
\(H(n)={}_nH(m)\); and}
\item[\SGARefTarget{sga3:xi:lemma:4:5:item:2}\(2^\circ\)]
\textit{\(H(n)_0={}_nH_0\).}
\end{enumerate}
\textit{Then there exists a unique multiplicative-type subgroup
\(H\) of \(G\) such that \({}_nH=H(n)\) for every \(n\).}

Uniqueness is already contained in \textit{\SGARefLink{sga3:xi:theorem:4:1:condition:b}{(b)}}.  For existence, an
immediate induction reduces us to the case in which
\(I\mathfrak m=0\), where \(\mathfrak m\) is the maximal ideal of
\(A\).\footnote{The printed source abruptly writes
\(J\mathfrak m=0\) here and uses \(J\) throughout the ensuing
argument, although only \(I\) was introduced in the statement.  We
retain \(I\), the same ideal, throughout.} Put
\[
  k=A/\mathfrak m,\qquad
  S_0=\operatorname{Spec}(k),\qquad
  G_0=G\times_SS_0,
\]
and let \(\mathfrak g_0\) and \(\mathfrak h_0\) be the Lie algebras
of \(G_0\) and \(H_0\), respectively.  There is a canonical
isomorphism of groups
\SGARefTarget{sga3:xi:display:4:5:infinitesimal-kernel}
\[
  \mathfrak g_0\otimes_kI
  \simeq
  \operatorname{Ker}\bigl(G(S)\longrightarrow G(S_0)\bigr);
\]
see Exposé III.  By \SGARefLink{sga3:ix:theorem:3:6:bis}{Exposé IX, 3.6 bis} and \SGARefLink{sga3:xi:proposition:3:7}{3.7}, there exists a
multiplicative-type subgroup \(H\) of \(G\) reducing to \(H_0\), and
such an \(H\) is determined, up to an inner automorphism by an
element of
\[
  N=\operatorname{Ker}\bigl(G(S)\longrightarrow G(S_0)\bigr).
\]
Thus the set \(P\) of liftings \(H\) of \(H_0\) is a principal
homogeneous set under \(N/M\), where
\[
  M=\{g\in N\mid \operatorname{int}(g)\mathbin{\cdot}H=H\}.
\]

One sees readily (Exposé III) that this subgroup is precisely the
vector subspace of \(\mathfrak g_0\otimes_kI\) consisting of the
invariants under \(H_0\), where \(H_0\) acts on
\(\mathfrak g_0\otimes_kI\) through the representation induced by
the adjoint representation of \(G_0\).  Likewise, the set \(P(n)\)
of liftings of \(H(n)_0={}_nH_0\) is a principal homogeneous set
under \(N/M(n)\), where \(M(n)\) is the vector subspace of
\(\mathfrak g_0\otimes_kI=N\) consisting of the invariants under
\({}_nH_0\).  The density theorem, \SGARefLink{sga3:ix:theorem:4:7}{Exposé IX, 4.7}, shows readily
that \(M=M(n)\) for sufficiently large \(n\), where ``large'' is
understood with respect to the divisibility order on the positive
integers.  Consequently, the natural map
\[
  H\longmapsto{}_nH
\]
from \(P\) to \(P(n)\), which is compatible with the actions of
\(N\), and hence with the homomorphism
\[
  N/M\longrightarrow N/M(n)
\]
between the acting groups, is bijective for sufficiently large
\(n\).  The conclusion of
\SGARefLink{sga3:xi:lemma:4:5}{4.5} follows immediately.

To verify conditions \textit{\SGARefLink{sga3:xi:proposition:3:7:item:c}{(c)}} and \textit{\SGARefLink{sga3:xi:proposition:3:7:item:d}{(d)}} of
\SGARefLink{sga3:xi:proposition:3:7}{3.7}, we use
\SGARefLink{sga3:xi:corollary:3:8}{3.8}, which reduces us to checking
\textit{(\(c'\))} and \textit{(\(d'\))} below.

\smallskip
\noindent\SGARefTarget{sga3:xi:theorem:4:1:condition:c-prime}
\textit{(\(c'\))} The \(T_n\) are smooth over \(S\), and the
transition morphisms \(p_{n,m}\colon T_m\to T_n\) are smooth.

This is precisely
\SGARefLink{sga3:xi:corollary:2:2:bis}{2.2 bis}.

\smallskip
\noindent\SGARefTarget{sga3:xi:theorem:4:1:condition:d-prime}
\textit{(\(d'\))} With the notation of
\SGARefLink{sga3:xi:corollary:3:8}{3.8}, a point \(\xi\) of \(F\)
with values in a field \(k\) over \(S\) is just a
multiplicative-type subgroup \(H_0\) of \(G_0=G_k\).  Repeating the
argument in \SGARefLink{sga3:xi:theorem:4:1:condition:b-prime}{\textit{(\(b'\))}} above, we see that the integer
\(d(\xi)\) considered in
\SGARefLink{sga3:xi:corollary:3:8}{3.8} is
\SGARefTarget{sga3:xi:display:4:1:d-xi-quotient}
\[
  d(\xi)=\dim_k\bigl(\mathfrak g_0/\mathfrak g_0^{H_0}\bigr),
\]
where \(\mathfrak g_0\) is the Lie algebra of \(G_0\) and
\(\mathfrak g_0^{H_0}\) is the vector subspace of invariants under
\(H_0\).  This is a finite integer, so condition
\SGARefLink{sga3:xi:corollary:3:8:item:d-prime:1}
{\(\textit{(d')}1^\circ\textit{)}\)} of
\SGARefLink{sga3:xi:corollary:3:8}{3.8} holds.

With the notation of
\SGARefLink{sga3:xi:corollary:3:8:item:d-prime:2}
{\(\textit{(d')}2^\circ\textit{)}\)}, giving a morphism \(v\colon
X\to F\) amounts to giving a multiplicative-type subgroup \(H\) of
\(G_X\).  For \(x\in X\), the integer \(d(\xi_x)\)\footnote{The printed source gives only
\(\bigl(\mathfrak g\otimes\kappa(x)\bigr)^{H_x}\) here and later
\(\operatorname{rank}_x\mathfrak g^H\), omitting the quotient.  That
reading contradicts both \SGARefLink{sga3:xi:display:4:1:d-xi-quotient}{the immediately preceding formula}
\(d(\xi)=\dim(\mathfrak g_0/\mathfrak g_0^{H_0})\) and the lifting
torsor \(N/M\).  The quotient is restored in both places.} is then
the dimension of
\SGARefTarget{sga3:xi:display:4:1:d-xi-x-invariants}
\[
  \bigl(\mathfrak g\otimes\kappa(x)\bigr)
  \big/
  \bigl(\mathfrak g\otimes\kappa(x)\bigr)^{H_x}.
\]
where \(\mathfrak g\) is the sheaf of Lie algebras of \(G_X\).  It
is locally free of finite type over \(X\), because \(G_X\) is smooth
over \(X\), and \(\mathfrak g\otimes\kappa(x)\) is the Lie algebra of
the fibre \(G_x\) of \(G_X\) at \(x\); here \(H_x\) is the fibre of
\(H\) at \(x\).  Now \(\mathfrak g\) is a locally free module on
\(X\)\footnote{The printed source says ``over \(S\)'' here, although
\(\mathfrak g\) is the Lie-algebra sheaf of \(G_X\); the base is
therefore \(X\).} on which the multiplicative-type group \(H\) acts, and the
invariant subfunctor \(\mathfrak g^H\) is therefore given by a
locally direct-summand, hence locally free, subsheaf of
\(\mathfrak g\).  By descent, one reduces to the case in which
\(H\) is diagonalizable, and then applies \SGARefLink{sga3:i:proposition:4-7-3}{Exposé I, 4.7.3}, noting
that the invariant subsheaf corresponds to the degree-zero
component.  Consequently,
\[
  d(\xi_x)=\operatorname{rank}_x
  \bigl(\mathfrak g/\mathfrak g^H\bigr),
\]
so this is a locally constant function of \(x\).  This proves
condition \SGARefLink{sga3:xi:theorem:4:1:condition:d-prime}{\textit{(\(d'\))}}.

We have thus verified the conditions of
\SGARefLink{sga3:xi:proposition:3:7}{3.7}, completing the proof of
\SGARefLink{sga3:xi:theorem:4:1}{4.1}.

\medskip
\noindent\SGARefTarget{sga3:xi:remark:4:6}\textbf{Remark 4.6.} ---
When \(G=\operatorname{GL}(n)_S\), one can give a much simpler and
more explicit direct proof of
\SGARefLink{sga3:xi:theorem:4:1}{4.1}, using \SGARefLink{sga3:i:proposition:4-7-3}{Exposé I, 4.7.3}.  The
proof moreover shows that, in this case, the moduli scheme is a
disjoint sum of affine \(S\)-schemes.  Proceeding as indicated in
\SGARefLink{sga3:xi:remark:4:3}{4.3}, one obtains the same result
whenever \(G\) is a closed subgroup of a group of the form
\(\operatorname{GL}(n)_S\).

One must not, however, suppose that the preschemes representing the
functors in \SGARefLink{sga3:xi:theorem:4:1}{4.1} or
\SGARefLink{sga3:xi:corollary:4:2}{4.2} are always disjoint sums of
affine \(S\)-schemes.  For example, let \(H\) be a
multiplicative-type group of finite type over a locally noetherian
prescheme \(S\).  By \SGARefLink{sga3:x:proposition:5-11}{Exposé X, 5.11}, \(H\) is isotrivial if and only
if
\[
  \operatorname{Hom}_{S\text{-}\mathrm{gr}}(H,\mathbb G_m)
\]
is a disjoint sum of preschemes affine over \(S\).  But, as noted in
\SGARefLink{sga3:x:paragraph:1:6}{Exposé X, 1.6}, there may exist nonisotrivial tori \(H\) of relative
dimension \(2\).  For such a torus,
\(\operatorname{Hom}_{S\text{-}\mathrm{gr}}(H,\mathbb G_m)\) is not
a disjoint sum of \(S\)-preschemes affine over \(S\), and likewise
the ``dual'' twisted constant group
\[
  \operatorname{Hom}_{S\text{-}\mathrm{gr}}(\mathbb G_m,H)
\]
is not such a sum.  Indeed, if two commutative twisted constant
groups \(R,R'\) of finite presentation are dual to one another,
\[
  R'=\operatorname{Hom}_{S\text{-}\mathrm{gr}}(R,\mathbb Z_S),
\]
then one readily sees that one is isotrivial if and only if the
other is.

This last point also shows that such an \(H\) is not isomorphic to a
subgroup of a group of the form \(\operatorname{GL}(n)_S\).  More
precisely, one can show that a multiplicative-type group of finite
type over a connected locally noetherian \(S\) is isomorphic to a
subgroup of a group of the form
\(\operatorname{Aut}_{\mathrm{Mod}}(E)\), where \(E\) is a locally
free module of finite type over \(S\), if and only if it is
isotrivial.  Finally, taking \(G=H\times\mathbb G_m\) in the two
preceding examples gives an example in which the prescheme
representing the functor \(F\) of
\SGARefLink{sga3:xi:theorem:4:1}{4.1} is not a disjoint sum of
affine \(S\)-preschemes; if desired, \(G\) may be taken to be a
torus of relative dimension \(3\).

